\documentclass[output=paper]{langscibook}
\ChapterDOI{10.5281/zenodo.22880234}
\title{Introduction}
\author{Ulrike Mosel}
\abstract{}
\begin{document}
\maketitle\label{chap:4}
The texts of this part were written between 2009 and 2012 for a school book about plants that for various reasons could not be completed, but text excerpts have been integrated in the Multifunctional Teop English Dictionary (\textsc{mted}) as ``examples''.\footnote{In \textcite{Mosel:2019bt}, search the Semantic Domain of the \textsc{mted} for ``plants'' and the Parts of Speech for ``\textsc{n}3'' and ``\textsc{n}3.\textsc{constr}'' and you will find 293 entries about plants and parts of plants.} The following 10 chapters document that in the beginning of the 21st century local plants are still an essential resource for the production of artefacts or are at least remembered as such by elder speakers.

In the first chapter Joyce Maion describes the growth of the coconut palm and the use of its parts, which confirms Shoffner's statement (\citeyear[284]{Shoffner:1976eg}) that the coconut palm ``has the most diverse uses of any member of the Teop plant assemblage. Almost every part of this tree can be utilized.'' 


The following chapters by Joyce Maion and Ruth Sima Rigamu deal with the sago palm, the Blackboard Tree, the young bread fruit tree, three kinds of bamboo, various plants that are used for making strings, and finally with three kinds of pandanus, see Table \ref{tab:plantsandartefacts}.

The bamboos described in Chapter \ref{chap:10} represent three distinct species called \textit{koopu}, \textit{koverau}, and \textit{kururu} with distinct usages. There is no generic term for them in the Teop language.


\begin{table}[t!]
    \caption{Plants and the artefacts produced from them}\label{tab:plantsandartefacts}
 \begin{tabularx}
{\textwidth}{llQ}

    \lsptoprule
        \textsc{chap}&\textsc{plants}&\textsc{artefacts}\\
    \midrule
        \ref{chap:5}&coconut palm&baskets, trays, brooms, strings, planting sticks\\
        \ref{chap:6}&sago palm&thatch roof\\
        \ref{chap:7}&blackboard tree&hull of dug-out canoes\\
        \ref{chap:8}&young breadfruit tree&strings for tying fishing nets\\
        \ref{chap:9}&sea poison tree&floats of fishing nets\\
        \ref{chap:10}&bamboos&spears, water containers, containers for cooking, walls of dwelling houses, and strips for sewing the thatch roof\\
        \ref{chap:11}&vines, aerial roots, fibres&strings for kite-fishing, strings for firmly lashing things together,  straps for backpacks, and mats\\
        \ref{chap:12}&hibiscus fibres&straps for backpacks, skirts for dancing\\
        \ref{chap:13}&pandanus&mats, baskets, hats\\
        \ref{chap:14}&giant pandanus&hoods\\
    \lspbottomrule
    \end{tabularx}
    \end{table}


Chapter \ref{chap:11} on ``Vines, aerial roots, fibres and strings'' deals in 14 subsections with plants, parts of plants, and strings that are formally classified as \textit{kara}, e.g. \textit{kara koris} `vine', \textit{kara kave} `net vine', \textit{kara ruhu} `aerial root of the Bengal Ficus', and \textit{kara obai} `string made of the \textit{obai} `Sweet Flag' (\textit{Acorus Calamus})'. The word \textit{kara} functions as the head of compounds and means `sth. that is thin and long'. Imported plastic strings are called \textit{kara sitiring} in contrast to \textit{kara naono} `string made of local plant materials'.\footnote{Search the \textsc{mted} Headword field for ``kara'' with a space after \textit{kara} and find 33 entries.}
\end{document}
